\documentclass[10pt,a4paper]{amsart}
\usepackage[margin=19mm]{geometry}
\usepackage{amsmath,amssymb,mathtools}
\usepackage[scr=rsfs,cal=euler]{mathalfa}
\usepackage{xfrac}
\usepackage[unicode,hidelinks,bookmarks=false]{hyperref}
\newcommand{\ie}{i.e.\ }
\newcommand{\cf}{cf.\ }
\newcommand{\resp}{respectively\ }
\newcommand{\Subsection}{\subsection}
\providecommand{\nobreakdash}{\nobreak\hbox{-}}
\setlength{\emergencystretch}{2em}
\multlinegap0pt
\usepackage{bm}
\usepackage{algorithm}\usepackage{algpseudocode}
\usepackage{float}
\usepackage{caption}
\usepackage{amssymb}
\usepackage[shortlabels]{enumitem}
\usepackage{multirow}
\newtheorem{lemma}{Lemma}
\title{On the Almost Sure Convergence of the Stochastic Three Points Algorithm}
\author{Taha EL BAKKALI EL KADI, Omar SAADI}
\date{Source: arXiv:2501.13886v5 (CC BY 4.0)}
\begin{document}
\maketitle

\noindent\emph{Source and licence.} On the Almost Sure Convergence of the Stochastic Three Points Algorithm. Version arXiv:2501.13886v5 (CC BY 4.0), distributed by arXiv under the Creative Commons Attribution 4.0 International licence, http://creativecommons.org/licenses/by/4.0/, as recorded in the arXiv licence metadata for this exact version and shown on the abstract page https://arxiv.org/abs/2501.13886v5. Full licence notice: \texttt{licenses/arxiv-2501.13886-CC-BY-4.0.txt}.\par\smallskip

\noindent\textit{Editorial scope.} Counted unit: the article's lemma Lem\_Reb (source0.tex lines 1216--1223). The article's complete original proof of it is delivered with it (source0.tex lines 1224--1259, one \texttt{proof} environment of 36 lines). No mathematics is written, repaired or re-derived: the statement and the proof are the article's own lines, in the article's own notation, and no retained line is substituted or re-flowed.  No further source-local context is needed: every cross-reference of every retained line resolves inside the two delivered spans.  Nothing was omitted from any retained span, so the package carries no omission record.  No retained line contains a citation, so the wrapper carries no bibliography and no bibliography entry.  No retained line contains a figure environment, an image-inclusion command or drawing code, so no image is delivered and the package declares no asset dependency.  Editorial formatting only, in addition: an AMS \texttt{amsart} preamble that restates the article's own declarations and macros used by the retained lines, namely the environments \texttt{lemma} on separate counters, together with the standard packages those lines need.  The retained environments print in this document as Lemma 1 in order of appearance, whereas the article prints the same items with the numbering of its own full document.  The complete native source of the article as distributed in the arXiv e-print for this exact version is bundled unchanged as \texttt{sources/172062/source0.tex}.
\begin{lemma}\label{Lem_Reb} 
 Let \( \{X_t\}_{t \ge 1} \) be a sequence of nonnegative real numbers that is non increasing and converges to \( 0 \), and let \( \{\alpha_t\}_{t \ge 1} \) be a sequence of real numbers such that  \( \sum_{t=1}^{\infty} \alpha_t X_t \) converges. Then, we have:
\[
X_T = o\left( \frac{1}{\sum_{t=1}^T \alpha_t} \right).
\]


\end{lemma}

\begin{proof}
    For all $T \ge 1$, we define $U_T=X_T \sum_{i=1}^T \alpha_i$ and $R_T=\sum_{i=T}^{\infty} \alpha_i X_i.$ We then have:
 
 $$
    U_T=X_T\sum_{i=1}^T (R_{i}-R_{i+1})\frac{1}{X_i} .
  $$
 Let $T \ge 2.$ We have:
 \begin{align*}
      U_T&=X_T \left[ \sum_{i=1}^T R_i\frac{1}{X_i}-\sum_{i=1}^T R_{i+1}\frac{1}{X_i}  \right]\\
      &= X_T \left[ \sum_{i=1}^T R_i\frac{1}{X_i}-\sum_{i=2}^{T+1} R_i\frac{1}{X_{i-1}}   \right]\\
      &= X_T \left[ R_1 \frac{1}{X_1}-\frac{R_{T+1}}{X_{T}} +\sum_{i=2}^T R_i (\frac{1}{X_i}-\frac{1}{X_{i-1}})  \right]\\
      &=R_1 \frac{X_T}{X_1}-R_{T+1}+X_T \sum_{i=2}^T R_i (\frac{1}{X_i}-\frac{1}{X_{i-1}}).
 \end{align*}
 To prove \( \lim_{T \to +\infty} U_T = 0 \), it suffices to show that:
\[
\lim_{T \to +\infty} X_T \sum_{i=2}^T R_i (\frac{1}{X_i}-\frac{1}{X_{i-1}}) = 0.
\]

Let $\epsilon>0$ and $T_0 \ge 2$ such that for all $T \ge T_0$, we have $ R_T  \le \frac{\epsilon}{2}$. Let $T > T_0$, we have:
 \begin{align*}
 \left| X_T \sum_{i=2}^T R_i (\frac{1}{X_i}-\frac{1}{X_{i-1}})    \right| &\le X_T \sum_{i=2}^{T_0} |R_i|(\frac{1}{X_i}-\frac{1}{X_{i-1}}) +X_T\sum_{i=T_0+1}^{T} \frac{\epsilon}{2} (\frac{1}{X_i}-\frac{1}{X_{i-1}})
 \\
 &= X_T \sum_{i=2}^{T_0} |R_i|(\frac{1}{X_i}-\frac{1}{X_{i-1}}) + \frac{\epsilon X_T}{2}  (\frac{1}{X_T}-\frac{1}{X_{T_0}})
  \\
 &= X_T \sum_{i=2}^{T_0} |R_i|(\frac{1}{X_i}-\frac{1}{X_{i-1}}) + \frac{\epsilon}{2} (1-\frac{X_T}{X_{T_0}})
 \\&\le X_T \sum_{i=2}^{T_0} |R_i|(\frac{1}{X_i}-\frac{1}{X_{i-1}}) + \frac{\epsilon}{2}.
 \end{align*}
As $\lim_{T \to +\infty} X_T=0$, there exists $T_1 \ge T_0$ such that for all $T \ge T_1$,
 \begin{align*}
 \left| X_T \sum_{i=2}^T R_i (\frac{1}{X_i}-\frac{1}{X_{i-1}})    \right| 
 &\le X_T \sum_{i=2}^{T_0} |R_i|(\frac{1}{X_i}-\frac{1}{X_{i-1}}) + \frac{\epsilon }{2}\le \epsilon.
 \end{align*}
 Therefore $\lim_{T \to +\infty} X_T \sum_{i=2}^T R_i (\frac{1}{X_i}-\frac{1}{X_{i-1}}) = 0,$ and we deduce that: \[
X_T = o\left( \frac{1}{\sum_{t=1}^T \alpha_t} \right) \text{ as } T \to +\infty.
\]
\end{proof}

\end{document}
